\section{投资问题：利润应由多少股分享}

恒逸的静态总股本不是正确的每股估值分母。回购库存股不参与当前经济权益，可转债在15.06元股价下深度价内并可能转成新股，第七期员工持股计划又可能让回购专户股份回流。核心判断是，主模型必须同时处理库存股扣除、转债潜在新增、员工计划上限及其现金流入；任何只计稀释或只计降债／现金的做法都不完整。

截至2026-06-30，公司法定总股本38.21562147亿股；分红登记口径回购库存股3.81840072亿股，得到经济股本34.39722075亿股。剩余“恒逸转2”本金29.944338亿元、转股价10.32元，对应潜在新增2.90158314亿股，转债股数敏感性达到37.29880389亿股。再加入已获批准的第七期员工计划上限1.508138亿股，得到主模型已披露稀释股本38.80694189亿股。

\begin{exhibitbox}[三层股本桥]
\noindent\begin{tabularx}{\textwidth}{L{4.3cm}R{2.3cm}X}
\toprule
\textbf{步骤} & \textbf{亿股} & \textbf{计算与用途} \\
\midrule
法定总股本 & 38.21562147 & 2026-06-30官方股本 \\
减：回购库存股 & -3.81840072 & 不参与当前经济权益 \\
第一层：经济股本 & 34.39722075 & 当前未转股、库存股扣除视角 \\
加：剩余转债潜在股份 & +2.90158314 & 29.944338亿元／10.32元 \\
第二层：转债股数敏感性 & 37.29880389 & 不回加未量化有效利息 \\
加：第七期员工计划上限 & +1.50813800 & 回购专户非交易过户上限 \\
第三层：已披露稀释股本 & 38.80694189 & 主估值分母 \\
\bottomrule
\end{tabularx}
\end{exhibitbox}
\sourcenote{HY-OFF-CB（股本、转债余额和10.32元转股价，截至2026-06-30）；HY-OFF-DIV（回购库存股登记口径）；HY-OFF-ESOP7（计划上限）。}

三层桥不是会计准则下精确的加权平均稀释EPS，而是面向估值的已披露股份敏感性。转股和员工计划实际过户具有时点差，历史EPS仍应采用公司会计口径；目标价则必须防止以更小分母夸大未来每股价值。

\section{转债：稀释与降债必须同时更新}

“恒逸转2”截至2026-06-30剩余本金约29.94亿元，转股价10.32元，转股期持续至2028-07-20。15.06元盘中价高于转股价，转股经济性较强。若全部转股，股数增加约2.90亿股，同时应付债券和未来利息负担下降；只把股数放进分母会低估降债价值，只回加利息又不加股数则高估EPS。

主模型采取保守且可复算的处理：加入潜在股数，但不回加任何转债利息。原因是财务报表中的利息不只有票面利率，还包括实际利率法下的折价摊销，税务影响也未单列。按票面利率简单回加会制造伪精确。转债实际转股后，应以已转股数量、债务减少和财务费用变化同时替换当前敏感性。

\begin{keyinsight}[转债的“所以呢”]
转股不是纯利空，也不是纯利好。\textbf{行动上，任何转股公告都要同时更新三项：新增股数、剩余债务和实际利息；在公司未量化有效利息前，主模型不作利润回加。}
\end{keyinsight}

\section{员工计划：最大稀释与最大现金是一组假设}

{\sloppy 第七期员工持股计划拟受让不超过1.508138亿股回购股份。受让价为12.43元／股，对应最高现金18.746155亿元。计划存续期不超过36个月，取得股票锁定期12个月；参加资金来自员工合法薪酬、自筹资金及法律允许方式。最终参加人、认购量、过户时点可能低于或晚于上限。\par}

主模型采用上限股份和相应最高现金，是为了避免单边处理：若把1.508138亿股全部加入分母，就应将同一假设下18.746155亿元现金加入权益价值；若最终只认购部分，应按实际股数和现金同步下调。该处理并不判断员工计划能否提升经营，也不把激励效果预先写入利润。

\begin{exhibitbox}[第七期员工计划的估值与会计边界]
\noindent\begin{tabularx}{\textwidth}{L{3.45cm}L{2.45cm}L{2.8cm}X}
\toprule
\textbf{字段} & \textbf{披露值} & \textbf{主模型处理} & \textbf{剩余风险} \\
\midrule
股份上限 & 1.508138亿股 & 全部加入稀释分母 & 最终认购可能不足 \\
受让价 & 12.43元／股 & 用于现金桥 & 过户时点与价格不再重估 \\
最高现金 & 18.746155亿元 & 加入权益价值 & 实际现金随认购变化 \\
存续／锁定 & 不超36个月／12个月 & 说明供给时点 & 锁定后处置节奏未知 \\
资金来源 & 员工薪酬、自筹等 & 不视为公司新增债务 & 员工缴款完成度未知 \\
股份支付 & 按准则处理 & 主模型不回加、不另造费用 & 2026费用金额和摊销未知 \\
\bottomrule
\end{tabularx}
\end{exhibitbox}
\sourcenote{HY-OFF-ESOP7（2026-06-13）；现金为1.508138亿股乘12.43元／股，模型截止2026-07-22。}

\begin{exhibitbox}[员工计划认购比例敏感性]
\small
\noindent\begin{tabularx}{\textwidth}{L{2.6cm}R{2.0cm}R{2.0cm}R{2.1cm}X}
\toprule
\textbf{认购比例} & \textbf{估值股数} & \textbf{匹配现金} & \textbf{基准价值} & \textbf{费用边界} \\
\midrule
0\% & 37.2988亿股 & 0 & 15.5836元 & 不发生过户则无该计划费用 \\
50\% & 38.0529亿股 & 9.3731亿元 & 15.5211元 & 股份支付仍待披露 \\
100\%（主模型） & 38.8069亿股 & 18.7462亿元 & 15.4611元 & 股份支付仍待披露 \\
\bottomrule
\end{tabularx}
\end{exhibitbox}
\sourcenote{HY-OFF-ESOP7；AStock House敏感性。每股价值公式为：（77.5×7.5＋匹配现金）÷估值股数；不预设股份支付费用。}

股份支付费用是本案必须明确的未知项。计划文件确认将按股份支付准则和公司会计政策处理，却没有给出授予日公允价值、归属／服务条件、费用总额和2026摊销。House 77.5亿元为报告口径利润，不将未知费用回加，也不虚构一个费用点估计；这部分缓冲包含在H2正常化中。

\section{剩余2.31亿库存股：用途未确认}

回购库存股3.81840072亿股扣除第七期计划上限1.508138亿股后，仍有2.31026272亿股用途未确认。它们当前仍从经济股本中扣除，但不能假设永久注销。未来可能继续用于员工计划、出售、注销或其他合规安排，具体取决于公司公告。

剩余库存股的风险来自选择权不对称：若未来回流，分母增加，通常也伴随现金、激励或资本运作价值；若注销，现有股东每股权益提高。没有披露前，主模型不把它们加入基准，也不提前赋予注销溢价，而是把2.31亿股列为边界敏感性。

\begin{exhibitbox}[剩余库存股情景]
\noindent\begin{tabularx}{\textwidth}{L{2.6cm}L{3.1cm}L{3.1cm}X}
\toprule
\textbf{未来用途} & \textbf{股本影响} & \textbf{价值补偿} & \textbf{当前处理} \\
\midrule
继续员工计划／激励 & 经济股本回流 & 现金与潜在激励效果 & 等正式方案，股数／现金／费用同时更新 \\
市场出售或其他安排 & 股份供给增加 & 可能取得现金 & 等公告和价格 \\
注销 & 每股权益上升 & 无现金流入 & 不提前给予注销信用 \\
继续库存 & 当前经济股本不变 & 无新增现金 & 维持边界风险 \\
\bottomrule
\end{tabularx}
\end{exhibitbox}
\sourcenote{HY-OFF-DIV；HY-OFF-ESOP7；analysis/valuation\_audit.md。2.31026272亿股为3.81840072减1.508138，截止2026-07-22用途未确认。}

\begin{exhibitbox}[剩余库存股回流的量化边界]
\small
\noindent\begin{tabularx}{\textwidth}{L{3.2cm}R{2.0cm}R{2.0cm}R{2.0cm}X}
\toprule
\textbf{情景} & \textbf{估值股数} & \textbf{新增现金} & \textbf{基准价值} & \textbf{费用／解释} \\
\midrule
继续库存或注销 & 38.8069亿股 & 0 & 15.4611元 & 当前经济分母不变 \\
无现金补偿全部回流 & 41.1172亿股 & 0 & 14.5923元 & 纯下行敏感性；费用未知 \\
按12.43元出售 & 41.1172亿股 & 28.7166亿元 & 15.2907元 & 未计交易税费 \\
按15.06元出售 & 41.1172亿股 & 34.7926亿元 & 15.4385元 & 未计交易税费 \\
\bottomrule
\end{tabularx}
\end{exhibitbox}
\sourcenote{AStock House敏感性，截止2026-07-22。通式为$[599.996155+2.31026272\times P]\div41.11720461$；$P$为每股回流现金价格。}

无现金补偿回流会把基准价值压至14.59元，足以吞没当前15.7元目标相对市价的有限空间；这正是剩余库存股必须作为正式边界而不能只作定性脚注的原因。若未来用于激励，还需把股份支付费用与股数、现金同时更新。

\section{每股价值的更新纪律}

主模型基准EPS为77.5亿元除以38.80694189亿股，约1.997元。员工计划现金加入情景权益价值，而不加入利润。转债有效利息和股份支付费用均不回加。这样的处理可能偏保守，却避免在信息不全时同时使用最小股本、最高利润和全部现金。

\begin{riskbox}[资本结构失效条件]
若剩余库存股回流但现金／激励价值不足、员工计划费用显著高于H2缓冲，或可转债大规模转股而债务和利息未同步下降，则每股价值低于当前模型。\textbf{行动上，收到正式公告后立即重算股本、现金和利润三张表，不等待年度报告。}
\end{riskbox}

相反，若转股显著降低净债务与利息、员工计划实际认购低于上限且股份支付费用温和、其余库存股注销，稀释折价可以收窄。但这些均不是当前目标价的前提。
